\chapter{The high-temperature cluster expansion}
\label{chap:cluster_expansion}

Hastings \cite[Section~IV]{Hastings2006Solving} approximates the thermal
operator $e^{-\beta H}$ of a Hamiltonian $H=\sum_k H_k$ by expanding it into
words in the terms $H_k$, joining the terms of a word into clusters when they
are adjacent in a graph on the terms, and discarding the words that contain a
large cluster. This chapter proves the parts of that expansion that involve no
sites: the cluster parts and their single-cluster bound, the truncated
expansion as an inclusion--exclusion sum over large clusters, the percolation
estimate for clusters of large diameter, the bound in trace norm, the error at
every temperature, and the expansion on a path of terms used for chains of
interaction labels \cite{Arad2016RigorousRG}. The terms are matrices indexed
by a finite set and the graph is arbitrary; for terms acting on sites it is
the overlap graph, in which two terms are adjacent when their supports meet.

\section{The single-cluster bound in Hastings' cluster expansion}
\label{sec:hastings_single_cluster_bound}



\begin{definition}[The cluster part]
    \label{def:hastings_cluster_part}
    \lean{ClusterExpansion.clusterPart,
        ClusterExpansion.clusterPart_eq_sum_powerset}
    \leanok
    For a Hamiltonian $H$ containing the terms $H_k$, a cluster $C$, its boundary
    $B(C)$ and $\beta\in\mathbb R$, put
    \begin{align}
        \rho(\beta,C)=\sum_{D\subseteq C}(-1)^{|D|}e^{-\beta(H-H_{B(C)}-H_D)}.
        \label{eq:hastings_cluster_part}
    \end{align}
    By inclusion and exclusion this is the sum of the terms of the power series
    of $e^{-\beta H}$ that contain every $H_i$ with $i\in C$ and no $H_j$ with
    $j\in B(C)$, the operator $\rho(\beta,C)$
    of~\cite[lines 1174--1177]{Hastings2006Solving}. With $E=C\setminus D$ it is
    $\sum_{E\subseteq C}(-1)^{|C|-|E|}e^{-(Y+\sum_{i\in E}\beta H_i)}$ for
    $Y=\beta(H-H_C-H_{B(C)})$.
\end{definition}

\begin{lemma}[Removing a cluster and its boundary]
    \label{thm:hastings_cluster_removal_trace}
    \lean{ClusterExpansion.re_trace_exp_sub_sum_sub_sum_le}
    \leanok
    Let $H$ be Hermitian, let the terms $H_k$ with $k\in C\cup B(C)$ be Hermitian
    with $\|H_k\|\le J$, and let $\beta\ge0$. Then, with
    $Z=\operatorname{tr}e^{-\beta H}$,
    \begin{align}
        \operatorname{tr}e^{-\beta(H-H_C-H_{B(C)})}\le e^{(n_C+n_{B(C)})J\beta}Z.
        \label{eq:hastings_cluster_removal_trace}
    \end{align}
\end{lemma}

\begin{proof}
    \leanok
    \uses{thm:trace_exp_add_le}
    The triangle inequality gives
    $\|\beta(H_C+H_{B(C)})\|\le(n_C+n_{B(C)})J\beta$, and the bound for a
    Hermitian perturbation in the exponent (Theorem~\ref{thm:trace_exp_add_le}),
    applied to $-\beta H+\beta(H_C+H_{B(C)})$, gives the claim.
\end{proof}

\begin{lemma}[The boundary count in the exponent]
    \label{thm:hastings_cluster_exp_count}
    \lean{ClusterExpansion.exp_card_add_card_mul_le}
    \leanok
    Let $\|H_k\|\le J$ for $k\in C\cup B(C)$, let $\beta\ge0$, and suppose that
    $n_C+n_{B(C)}\le Sn_C$. Then
    $e^{(n_C+n_{B(C)})J\beta}\le\bigl(e^{SJ\beta}\bigr)^{n_C}$.
\end{lemma}

\begin{proof}
    \leanok
    If $C\cup B(C)$ is empty, both sides equal $1$. Otherwise $J\ge0$, so
    $(n_C+n_{B(C)})J\beta\le Sn_CJ\beta$.
\end{proof}

\begin{lemma}[The bound on the cluster part]
    \label{thm:hastings_cluster_part_bound}
    \lean{ClusterExpansion.norm_trace_mul_clusterPart_le}
    \leanok
    \uses{def:hastings_cluster_part}
    Let $H$ be Hermitian, let the terms $H_k$ with $k\in C\cup B(C)$ be Hermitian
    with $\|H_k\|\le J$, and let $\beta\ge0$. Then for every operator $O$,
    with $Z=\operatorname{tr}e^{-\beta H}$,
    \begin{align}
        \left|\operatorname{tr}\bigl(O\rho(\beta,C)\bigr)\right|
        \le\|O\|\bigl(e^{\beta J}-1\bigr)^{n_C}e^{(n_C+n_{B(C)})J\beta}Z.
        \label{eq:hastings_cluster_part_bound}
    \end{align}
    This is the bound stated before equation (percb2)
    in~\cite[lines 1174--1195]{Hastings2006Solving}, with $\|O\|Z$ in place of the
    printed $\operatorname{tr}(Oe^{-\beta H})$, which is the form used in (percb2),
    and with $\|H_C+H_{B(C)}\|\le(n_C+n_{B(C)})J$. The printed form fails: applied
    to $\pm O$ it would make $\rho(\beta,C)$ a multiple of $e^{-\beta H}$
    \cite{gap:hastings05_cluster_part_trace_pairing}.
\end{lemma}

\begin{proof}
    \leanok
    \uses{def:hastings_cluster_part, thm:hastings_cluster_removal_trace, thm:trace_exp_alternating_sum}
    By~(\ref{eq:hastings_cluster_part}) with $E=C\setminus D$, the alternating-sum
    bound of Theorem~\ref{thm:trace_exp_alternating_sum},
    applied with $Y=\beta(H-H_C-H_{B(C)})$ and the terms $\beta H_i$, $i\in C$, of
    norm at most $\beta J$, gives
    \begin{align}
        \left|\operatorname{tr}\bigl(O\rho(\beta,C)\bigr)\right|
        \le\|O\|\bigl(e^{\beta J}-1\bigr)^{n_C}\operatorname{tr}e^{-Y}.\notag
    \end{align}
    Lemma~\ref{thm:hastings_cluster_removal_trace} gives
    $\operatorname{tr}e^{-Y}\le e^{(n_C+n_{B(C)})J\beta}Z$.
\end{proof}

\begin{lemma}[The boundary count]
    \label{thm:hastings_cluster_boundary_count}
    \lean{ClusterExpansion.clusterBoundary,
        ClusterExpansion.card_add_card_clusterBoundary_le}
    \leanok
    Let $C$ be a set of sites, and let $B(C)$ be the set of $j\notin C$ such
    that $H_j$ acts on some site of $C$. Suppose that every site $i\in C$ has a
    set $N(i)$ of at most $S$ sites that contains $i$ and every $j$ for which
    $H_j$ acts on $i$. Then
    \begin{align}
        n_C+n_{B(C)}\le Sn_C.\notag
    \end{align}
    In~\cite[lines 181--183]{Hastings2006Solving}, $N(i)$ is the set of sites
    within the interaction range $R$ of $i$, of which there are at most $S$,
    and $H_j$ acts only on sites within distance $R$ of $j$. This is the second
    inequality of~\cite[line 1186]{Hastings2006Solving}. The same count holds
    with $H_j$ acting on $i$ replaced by any relation between terms. For the
    boundary of terms whose support meets that of $H_C$, which the truncated
    expansion needs \cite{gap:hastings05_cluster_boundary}, $N(i)$ is the set
    of terms whose support meets that of $H_i$. Their centres lie within $2R$
    of $i$, so $S$ counts the sites within distance $2R$.
\end{lemma}

\begin{proof}
    \leanok
    The sets $C$ and $B(C)$ are disjoint, and both lie in
    $\bigcup_{i\in C}N(i)$, which has at most $Sn_C$ elements.
\end{proof}

\begin{lemma}[Equation (percb2)]
    \label{thm:hastings_cluster_part_percb2}
    \lean{ClusterExpansion.norm_trace_mul_clusterPart_le_pow}
    \leanok
    \uses{thm:hastings_cluster_part_bound}
    In the setting of Lemma~\ref{thm:hastings_cluster_part_bound}, suppose that
    $n_C+n_{B(C)}\le Sn_C$. Then, with $y=(e^{J\beta}-1)e^{SJ\beta}$,
    \begin{align}
        \left|\operatorname{tr}\bigl(O\rho(\beta,C)\bigr)\right|
        \le\|O\|y^{n_C}Z.
        \label{eq:hastings_cluster_part_percb2}
    \end{align}
    This is equation (percb2) of~\cite[lines 1186--1195]{Hastings2006Solving},
    multiplied by $Z$. The source obtains $n_C+n_{B(C)}\le Sn_C$ from
    the lattice, where $S$ bounds the number of terms acting on a site; here it
    is a hypothesis, which Lemma~\ref{thm:hastings_cluster_boundary_count}
    supplies.
\end{lemma}

\begin{proof}
    \leanok
    \uses{thm:hastings_cluster_part_bound, thm:hastings_cluster_exp_count}
    By Lemma~\ref{thm:hastings_cluster_exp_count},
    $e^{(n_C+n_{B(C)})J\beta}\le e^{SJ\beta n_C}$, and
    (\ref{eq:hastings_cluster_part_bound}) gives the claim.
\end{proof}

\begin{lemma}[Several clusters]
    \label{thm:hastings_cluster_part_several}
    \lean{ClusterExpansion.norm_trace_mul_clusterPart_biUnion_le}
    \leanok
    \uses{thm:hastings_cluster_part_percb2}
    Let $C_1,\dots,C_m$ be pairwise disjoint clusters with boundaries
    $B(C_1),\dots,B(C_m)$, each with $n_{C_i}+n_{B(C_i)}\le Sn_{C_i}$, and let
    $\rho(\beta,C_1,\dots,C_m)$ be the operator~(\ref{eq:hastings_cluster_part})
    for the cluster $C_1\cup\dots\cup C_m$ and the boundary
    $B(C_1)\cup\dots\cup B(C_m)$. Under the hypotheses of
    Lemma~\ref{thm:hastings_cluster_part_bound} on these terms,
    \begin{align}
        \left|\operatorname{tr}\bigl(O\rho(\beta,C_1,\dots,C_m)\bigr)\right|
        \le\|O\|\prod_{i=1}^m y^{n_{C_i}}Z.\notag
    \end{align}
    This is the bound on $\rho(\beta,C_1,C_2,\dots)$ that the remark
    of~\cite[lines 1218--1224]{Hastings2006Solving} derives in the same way as
    (percb2). For the clusters of one term of the power series, no $C_j$ meets
    $B(C_i)$, so the operator above is the source's $\rho(\beta,C_1,\dots,C_m)$.
\end{lemma}

\begin{proof}
    \leanok
    \uses{thm:hastings_cluster_part_percb2}
    The union of the clusters has $\sum_in_{C_i}$ sites, and the union of the
    boundaries has at most $\sum_in_{B(C_i)}$ terms, so the union satisfies
    $n_C+n_{B(C)}\le Sn_C$. Lemma~\ref{thm:hastings_cluster_part_percb2} gives
    $\|O\|y^{\sum_in_{C_i}}Z=\|O\|\prod_iy^{n_{C_i}}Z$.
\end{proof}

\begin{lemma}[The cluster part between thermal factors]
    \label{thm:hastings_cluster_part_thermal}
    \lean{ClusterExpansion.norm_trace_mul_exp_mul_clusterPart_mul_exp_le_pow}
    \leanok
    \uses{thm:hastings_cluster_part_percb2}
    In the setting of Lemma~\ref{thm:hastings_cluster_part_percb2}, let $n\ge1$ and
    $0\le k\le n-1$, and put $y=(e^{J\beta/n}-1)e^{SJ\beta/n}$. Then
    \begin{align}
        \left|\operatorname{tr}\bigl(Oe^{-k\beta H/n}\rho(\beta/n,C)
        e^{-(n-k-1)\beta H/n}\bigr)\right|
        \le\|O\|y^{n_C}Z.
        \label{eq:hastings_cluster_part_thermal}
    \end{align}
    This is the bound of~\cite[lines 1322--1333]{Hastings2006Solving}, multiplied by $Z$,
    where $\rho(\beta)=e^{-\beta H}$; the source proves it as (percb2).
\end{lemma}

\begin{proof}
    \leanok
    \uses{def:hastings_cluster_part, thm:hastings_cluster_removal_trace,
        thm:hastings_cluster_exp_count, thm:trace_exp_alternating_sum_between_pow}
    Put $Y=\beta(H-H_C-H_{B(C)})$. As in the proof of
    Lemma~\ref{thm:hastings_cluster_part_bound}, $\rho(\beta/n,C)$ is the alternating sum of
    the $e^{-(Y+\sum_{i\in E}\beta H_i)/n}$ over $E\subseteq C$. The bound for alternating
    sums between thermal factors (Theorem~\ref{thm:trace_exp_alternating_sum_between_pow}),
    with the weights $k/n$, $1/n$ and $(n-k-1)/n$ and the terms $\beta H_i$ of norm at most
    $\beta J$, gives
    \begin{align}
        \left|\operatorname{tr}\bigl(Oe^{-k\beta H/n}\rho(\beta/n,C)
        e^{-(n-k-1)\beta H/n}\bigr)\right|
        \le\|O\|\bigl(e^{J\beta/n}-1\bigr)^{n_C}Z^{(n-1)/n}
        \bigl(\operatorname{tr}e^{-Y}\bigr)^{1/n}.\notag
    \end{align}
    Lemma~\ref{thm:hastings_cluster_removal_trace} gives
    $\operatorname{tr}e^{-Y}\le e^{(n_C+n_{B(C)})J\beta}Z$. Taking
    $1/n$-th powers and using $Z^{(n-1)/n}Z^{1/n}=Z$ bounds the right side by
    $\|O\|(e^{J\beta/n}-1)^{n_C}e^{(n_C+n_{B(C)})J\beta/n}Z$, and
    Lemma~\ref{thm:hastings_cluster_exp_count}, with $\beta/n$ in place of
    $\beta$, bounds $e^{(n_C+n_{B(C)})J\beta/n}$ by $e^{SJ\beta n_C/n}$.
\end{proof}

\begin{lemma}[Products of cluster parts]
    \label{thm:hastings_cluster_part_product}
    \lean{ClusterExpansion.clusterPart_empty_empty,
        ClusterExpansion.norm_trace_mul_list_prod_clusterPart_le_pow}
    \leanok
    \uses{def:hastings_cluster_part}
    Let $H$ be a Hermitian Hamiltonian, let $\beta\ge0$, and put
    $Z=\operatorname{tr}e^{-\beta H}$. Let $n\ge1$, and for each $j<n$ let $C_j$ be a
    cluster with boundary $B_j$ such that the terms $H_k$ with $k\in C_j\cup B_j$ are
    Hermitian with $\|H_k\|\le J$, and $n_{C_j}+n_{B_j}\le Sn_{C_j}$. With
    $y=(e^{J\beta/n}-1)e^{SJ\beta/n}$, for every operator $O$,
    \begin{align}
        \left|\operatorname{tr}\bigl(O\rho(\beta/n,C_0)\rho(\beta/n,C_1)\cdots
        \rho(\beta/n,C_{n-1})\bigr)\right|
        \le\|O\|y^{\sum_jn_{C_j}}Z.
        \label{eq:hastings_cluster_part_product}
    \end{align}
    A factor with $C_j=B_j=\emptyset$ is the thermal factor $e^{-\beta H/n}$, so this bounds
    every product of thermal factors and cluster parts in the operator $\rho_m$
    of~\cite[lines 1334--1341]{Hastings2006Solving}.
\end{lemma}

\begin{proof}
    \leanok
    \uses{def:hastings_cluster_part, thm:hastings_cluster_removal_trace,
        thm:hastings_cluster_exp_count, thm:trace_exp_alternating_sum_product_pow}
    Put $Y_j=\beta(H-H_{C_j}-H_{B_j})$, so that $\rho(\beta/n,C_j)$ is the alternating sum of
    the $e^{-(Y_j+\sum_{i\in E}\beta H_i)/n}$ over $E\subseteq C_j$. The bound for products of
    alternating sums (Theorem~\ref{thm:trace_exp_alternating_sum_product_pow}), with
    the operator $O$ before the first factor, the identity before the others, the weight $1/n$
    on every factor and the terms $\beta H_i$ of norm at most $\beta J$, gives
    \begin{align}
        \left|\operatorname{tr}\bigl(O\rho(\beta/n,C_0)\cdots\rho(\beta/n,C_{n-1})\bigr)\right|
        \le\|O\|\prod_j\bigl(e^{J\beta/n}-1\bigr)^{n_{C_j}}
        \bigl(\operatorname{tr}e^{-Y_j}\bigr)^{1/n}.\notag
    \end{align}
    Lemmas~\ref{thm:hastings_cluster_removal_trace} and~\ref{thm:hastings_cluster_exp_count},
    the latter with $\beta/n$ in place of $\beta$, give
    $(\operatorname{tr}e^{-Y_j})^{1/n}\le e^{Sn_{C_j}J\beta/n}Z^{1/n}$, and the $n$ factors
    $Z^{1/n}$ multiply to $Z$.
\end{proof}

\section{The truncated expansion as an inclusion--exclusion sum}
\label{sec:hastings_cluster_truncation}

% Hastings, arXiv:cond-mat/0508554v3, proj3new.tex, lines 1090--1138 (active
% bonds, clusters, the boundary) and 1197--1225 (the inclusion--exclusion sum
% and the remark added).

Expanding $e^{-\beta H}=\sum_n\frac{(-\beta)^n}{n!}\bigl(\sum_kH_k\bigr)^n$ into
words $H_{k_1}\cdots H_{k_n}$, the truncated operator
$\rho(\beta,l_{\mathrm{proj}})$ keeps words containing no cluster of active
bonds of diameter larger than $l_{\mathrm{proj}}$
\cite[lines 1090--1114]{Hastings2006Solving} and equals the inclusion--exclusion
sum $\rho-\rho_1+\rho_2-\cdots$ over families of large clusters
\cite[lines 1218--1224]{Hastings2006Solving}, provided the cluster boundary
consists of terms whose supports meet the cluster support
\cite[lines 1131--1134]{Hastings2006Solving},
\cite{gap:hastings05_cluster_boundary}.

\begin{definition}[Clusters of a set of terms]
    \label{def:hastings_cluster_of_terms}
    \lean{SimpleGraph.overlapGraph,
        ClusterExpansion.boundary,
        ClusterExpansion.mem_boundary_overlapGraph,
        ClusterExpansion.IsCluster,
        ClusterExpansion.IsAdmissible}
    \leanok
    \uses{thm:hastings_cluster_boundary_count}
    Let $G$ be a graph on a finite term index set $K$ (in the overlap graph,
    $k\sim k'$ when $k\ne k'$ and their supports meet). For $C\subseteq K$, its
    \emph{boundary} is $\partial C=\{k\in K\setminus C:\exists c\in C,\ k\sim c\}$.
    A \emph{cluster} of $A\subseteq K$ is a $G$-connected subset $C\subseteq A$
    with $A\cap\partial C=\emptyset$. A family $U$ of subsets of $K$ is
    \emph{admissible} if $C\cap\partial C'=\emptyset$ for all $C,C'\in U$.
\end{definition}

\begin{lemma}[Clusters that meet]
    \label{thm:hastings_clusters_meet}
    \lean{ClusterExpansion.subset_of_connected_of_disjoint_boundary,
        ClusterExpansion.eq_of_isCluster_of_inter_nonempty,
        ClusterExpansion.pairwiseDisjoint_of_admissible,
        ClusterExpansion.biUnion_boundary_eq_of_admissible,
        ClusterExpansion.card_biUnion_of_admissible,
        ClusterExpansion.card_add_card_biUnion_boundary_le}
    \leanok
    \uses{def:hastings_cluster_of_terms}
    Intersecting clusters of $A$ coincide. For an admissible family $U$ of
    connected sets with $C_U=\bigcup_{C\in U}C$, the members of $U$ are pairwise
    disjoint, $\partial C_U=\bigcup_{C\in U}\partial C$, and if
    $n_C+n_{\partial C}\le Sn_C$ for all $C\in U$, then
    $n_{C_U}+n_{\partial C_U}\le Sn_{C_U}$.
\end{lemma}
\begin{proof}
    \leanok
    If $C$ is connected, $C\cap C'\ne\emptyset$, and $C\cap\partial C'=\emptyset$,
    any path in $C$ from $C\cap C'$ stays in $C'$, so $C\subseteq C'$; symmetry
    gives equality for intersecting clusters of $A$ and pairwise disjointness for
    admissible families, whence $\partial C_U=\bigcup_{C\in U}\partial C$ and
    the additive size bound.
\end{proof}

\begin{lemma}[Words containing a set of terms]
    \label{thm:hastings_words_inclusion_exclusion}
    \lean{ClusterExpansion.sum_pow_eq_sum_words,
        ClusterExpansion.sum_powerset_neg_one_pow_card_disjoint,
        ClusterExpansion.sum_words_contain_avoid}
    \leanok
    In any ring, $(\sum_{k\in s}a_k)^n=\sum_{w\in s^n}a_{w_1}\cdots a_{w_n}$,
    and for disjoint $B,C\subseteq K$, the sum of $a_{w_1}\cdots a_{w_n}$ over
    words of length $n$ whose letter set $X$ satisfies $C\subseteq X$ and
    $X\cap B=\emptyset$ equals
    \begin{align}
        \sum_{D\subseteq C}(-1)^{|D|}\Bigl(\sum_{k\notin B\cup D}a_k\Bigr)^n.
        \label{eq:hastings_words_inclusion_exclusion}
    \end{align}
\end{lemma}
\begin{proof}
    \leanok
    For $X\cap B=\emptyset$,
    $\sum_{D\subseteq C\setminus X}(-1)^{|D|}=[C\subseteq X]$.
\end{proof}

\begin{lemma}[The cluster part as a sum of words]
    \label{thm:hastings_cluster_part_words}
    \lean{ClusterExpansion.hasSum_div_factorial_smul_pow,
        ClusterExpansion.hasSum_words_clusterPart}
    \leanok
    \uses{thm:hastings_words_inclusion_exclusion, def:hastings_cluster_part}
    For $H=\sum_{k\in K}H_k$ and disjoint $B,C\subseteq K$, the words whose
    letter set contains $C$ and avoids $B$, weighted by $(-\beta)^n/n!$, sum to
    $\rho(\beta,C)$ of (\ref{eq:hastings_cluster_part}) with boundary $B$.
\end{lemma}
\begin{proof}
    \leanok
    Sum (\ref{eq:hastings_words_inclusion_exclusion}) weighted by
    $(-\beta)^n/n!$ over $n\ge0$.
\end{proof}

\begin{definition}[The truncated operator]
    \label{def:hastings_truncated_expansion}
    \lean{ClusterExpansion.NoLargeCluster,
        ClusterExpansion.largeConnected,
        ClusterExpansion.truncatedExpTerm,
        ClusterExpansion.truncatedExp}
    \leanok
    \uses{def:hastings_cluster_of_terms}
    For a set $T$ of large connected subsets of $K$, $\rho(\beta,l)$ is the sum
    over $n\ge0$ of $\frac{(-\beta)^n}{n!}\sum H_{k_1}\cdots H_{k_n}$ restricted
    to words whose letter set has no cluster in $T$
    \cite[lines 1101--1104]{Hastings2006Solving}.
\end{definition}

\begin{lemma}[Inclusion and exclusion for one word]
    \label{thm:hastings_word_large_clusters}
    \lean{ClusterExpansion.mem_largeConnected,
        ClusterExpansion.forall_isCluster_iff,
        ClusterExpansion.ite_noLargeCluster_eq_sum}
    \leanok
    \uses{thm:hastings_clusters_meet, def:hastings_truncated_expansion}
    For $A\subseteq K$, summing over admissible $U\subseteq T$ with
    $\bigcup_{C\in U}C\subseteq A$ and
    $A\cap\bigcup_{C\in U}\partial C=\emptyset$ gives
    \begin{align}
        [\text{no cluster of }A\text{ lies in }T]
        =\sum_{U}(-1)^{|U|}.
        \label{eq:hastings_word_large_clusters}
    \end{align}
\end{lemma}
\begin{proof}
    \leanok
    Such $U$ are precisely the subsets of the set $L$ of large clusters of $A$,
    and $\sum_{U\subseteq L}(-1)^{|U|}=[L=\emptyset]$.
\end{proof}

\begin{lemma}[The truncated operator by inclusion and exclusion]
    \label{thm:hastings_truncation_inclusion_exclusion}
    \lean{ClusterExpansion.hasSum_truncatedExpTerm,
        ClusterExpansion.truncatedExp_eq_sum,
        ClusterExpansion.exp_sub_truncatedExp}
    \leanok
    \uses{thm:hastings_word_large_clusters, thm:hastings_cluster_part_words,
        def:hastings_truncated_expansion}
    For $H=\sum_{k\in K}H_k$, the series $\rho(\beta,l)$ converges and, summing
    over admissible $U\subseteq T$ with $\rho(\beta,U)$ given by
    (\ref{eq:hastings_cluster_part}) for $\bigl(\bigcup_{C\in U}C,\bigcup_{C\in U}\partial C\bigr)$,
    \begin{align}
        \rho(\beta,l)=\sum_U(-1)^{|U|}\rho(\beta,U),\qquad
        e^{-\beta H}-\rho(\beta,l)=-\sum_{U\neq\emptyset}(-1)^{|U|}\rho(\beta,U)
        \label{eq:hastings_truncation_inclusion_exclusion}
    \end{align}
    \cite[lines 1197--1225]{Hastings2006Solving},
    \cite{gap:hastings05_cluster_boundary}.
\end{lemma}
\begin{proof}
    \leanok
    Substitute (\ref{eq:hastings_word_large_clusters}) into
    Definition~\ref{def:hastings_truncated_expansion}, exchange sums, and apply
    Lemma~\ref{thm:hastings_cluster_part_words} (with $\rho(\beta,\emptyset)=e^{-\beta H}$).
\end{proof}

\begin{lemma}[The bound summed by the percolation estimate]
    \label{thm:hastings_truncation_trace_bound}
    \lean{ClusterExpansion.norm_trace_mul_exp_sub_truncatedExp_le}
    \leanok
    \uses{thm:hastings_truncation_inclusion_exclusion,
        thm:hastings_cluster_part_several}
    If $H=\sum_kH_k$ has Hermitian terms with $\|H_k\|\le J$, $\beta\ge0$, and
    $n_C+n_{\partial C}\le Sn_C$ for all $C\in T$, then for every $O$, with
    $y=(e^{J\beta}-1)e^{SJ\beta}$ and $Z=\operatorname{tr}e^{-\beta H}$,
    \begin{align}
        \bigl|\operatorname{tr}\bigl(O(e^{-\beta H}-\rho(\beta,l))\bigr)\bigr|
        \le\|O\|Z\sum_{U\neq\emptyset}\prod_{C\in U}y^{n_C}.
        \label{eq:hastings_truncation_trace_bound}
    \end{align}
\end{lemma}
\begin{proof}
    \leanok
    \uses{thm:hastings_truncation_inclusion_exclusion,
        thm:hastings_cluster_part_several, thm:hastings_clusters_meet}
    Apply Lemmas~\ref{thm:hastings_cluster_part_several}
    and~\ref{thm:hastings_clusters_meet} to each term of
    (\ref{eq:hastings_truncation_inclusion_exclusion}).
\end{proof}

\begin{definition}[The operator of a set of terms]
    \label{def:hastings_cluster_operator}
    \lean{ClusterExpansion.clusterOp,
        ClusterExpansion.clusterOp_eq_clusterPart,
        ClusterExpansion.hasSum_words_clusterOp}
    \leanok
    \uses{thm:hastings_cluster_part_words}
    For $A\subseteq K$ and $H_D=\sum_{k\in D}H_k$, define
    \begin{align}
        E_A=\sum_{D\subseteq A}(-1)^{|A\setminus D|}e^{-\beta H_D},
        \label{eq:hastings_cluster_operator}
    \end{align}
    which sums the words in the expansion of $e^{-\beta H}$ whose letter set is
    exactly $A$ \cite[lines 1268--1269]{Hastings2006Solving}.
\end{definition}

\begin{lemma}[The truncated operator as a sum of products over clusters]
    \label{thm:hastings_truncation_cluster_products}
    \lean{ClusterExpansion.truncatedExp_eq_sum_clusterOp,
        ClusterExpansion.clusterOp_union,
        ClusterExpansion.clusters,
        ClusterExpansion.biUnion_clusters,
        ClusterExpansion.clusterOp_eq_noncommProd_clusters,
        ClusterExpansion.truncatedExp_eq_sum_noncommProd_clusters}
    \leanok
    \uses{def:hastings_cluster_operator, def:hastings_truncated_expansion,
        def:hastings_cluster_of_terms}
    Summing over $A\subseteq K$ with no cluster in $T$,
    $\rho(\beta,l)=\sum_AE_A$. If non-adjacent terms commute, then
    $E_A=\prod_{C\in\mathrm{clusters}(A)}E_C$ with commuting factors
    \cite[lines 1108--1110]{Hastings2006Solving}.
\end{lemma}
\begin{proof}
    \leanok
    \uses{def:hastings_cluster_operator, thm:hastings_clusters_meet}
    Grouping words by letter set $A$ gives $\rho(\beta,l)=\sum_AE_A$. For
    disjoint commuting sets $A,B$,
    $e^{-\beta H_{D_1\cup D_2}}=e^{-\beta H_{D_1}}e^{-\beta H_{D_2}}$ yields
    $E_{A\cup B}=E_AE_B$; induct over the clusters partitioning $A$.
\end{proof}

\section{The percolation sum in Hastings' cluster expansion}
\label{sec:hastings_percolation}

% Hastings, arXiv:cond-mat/0508554v3, proj3new.tex, lines 1196--1254 (the
% inclusion--exclusion sum over large clusters and eq. perc2).

At high temperature, \cite[lines 1196--1225]{Hastings2006Solving} bounds
$e^{-\beta H}-\rho(\beta,l_{\rm proj})$ via an inclusion--exclusion sum over
families $U$ of clusters of diameter $>l_{\rm proj}$, where each $U$ contributes
at most $\|O\|\prod_{C\in U}a_C$ with $0\le a_C\le y^{n_C}$ and
$y=(e^{J\beta}-1)e^{SJ\beta}$. For terms of range $R$, a cluster of $n$ terms
has diameter at most $2Rn$, so diameter $>l_{\rm proj}$ implies
$n_C\ge l+1$ with $l=\lfloor l_{\rm proj}/(2R)\rfloor$, correcting the
nearest-neighbour threshold $l_{\rm proj}+1$ of
\cite[line 1240]{Hastings2006Solving}; see
\cite{gap:hastings05_percolation_cluster_size}.

\begin{lemma}[Families of a given number of clusters]
    \label{thm:hastings_percolation_families}
    \lean{ClusterExpansion.sum_powersetCard_prod_le}
    \leanok
    For a finite set $T$ with weights $a_C\ge0$ ($C\in T$) and $m\in\mathbb N$,
    \begin{align}
        \sum_{U\subseteq T,\ |U|=m}\ \prod_{C\in U}a_C
        \le\frac1{m!}\Bigl(\sum_{C\in T}a_C\Bigr)^m
        \label{eq:hastings_percolation_families}
    \end{align}
    \cite[line 1210]{Hastings2006Solving}.
\end{lemma}
\begin{proof}
    \leanok
    Induct on $|T|$: for $b\notin T$ and $s=\sum_{C\in T}a_C$,
    $s^{m+1}/(m+1)!+a_bs^m/m!\le(s+a_b)^{m+1}/(m+1)!$.
\end{proof}

\begin{lemma}[All nonempty families of clusters]
    \label{thm:hastings_percolation_exp}
    \lean{ClusterExpansion.sum_nonempty_powerset_prod_eq,
        ClusterExpansion.sum_nonempty_powerset_prod_le_exp}
    \leanok
    For a finite set $T$ with weights $a_C\ge0$ ($C\in T$),
    \begin{align}
        \sum_{\emptyset\ne U\subseteq T}\ \prod_{C\in U}a_C
        =\prod_{C\in T}(1+a_C)-1
        \le\exp\Bigl(\sum_{C\in T}a_C\Bigr)-1
        \label{eq:hastings_percolation_exp}
    \end{align}
    \cite[lines 1238--1241]{Hastings2006Solving}.
\end{lemma}
\begin{proof}
    \leanok
    Expand $\prod_{C\in T}(1+a_C)$ and use $1+a_C\le e^{a_C}$.
\end{proof}

\begin{lemma}[The sum over large clusters]
    \label{thm:hastings_percolation_cluster_sum}
    \lean{ClusterExpansion.sum_pow_le_of_forall_lt,
        ClusterExpansion.sum_pow_le_of_card_le}
    \leanok
    If every $C\in T$ has $n_C\ge l+1$, $|\{C\in T:n_C=k\}|\le Nx^k$ for all
    $k$ ($N\ge0$), and $x,y\ge0$ with $xy<1$, then
    \begin{align}
        \sum_{C\in T}y^{n_C}
        \le N\frac{(xy)^{l+1}}{1-xy}
        \label{eq:hastings_percolation_cluster_sum}
    \end{align}
    \cite[lines 1238--1241]{Hastings2006Solving}.
\end{lemma}
\begin{proof}
    \leanok
    Group by $k=n_C\ge l+1$ and bound
    \begin{align}
        \sum_{k\ge l+1}N(xy)^k=N\frac{(xy)^{l+1}}{1-xy}.
        \notag
    \end{align}
\end{proof}

\begin{lemma}[The percolation bound]
    \label{thm:hastings_percolation}
    \lean{ClusterExpansion.sum_nonempty_powerset_prod_le_percolation,
        ClusterExpansion.pow_succ_div_one_sub_le,
        ClusterExpansion.pow_eq_exp_neg_div}
    \leanok
    \uses{thm:hastings_percolation_exp, thm:hastings_percolation_cluster_sum}
    In the setting of Lemma~\ref{thm:hastings_percolation_cluster_sum}, if
    $0\le a_C\le y^{n_C}$, $0<xy\le1/2$, and $\xi=1/\ln(1/(xy))$, then
    \begin{align}
        \sum_{\emptyset\ne U\subseteq T}\ \prod_{C\in U}a_C
        \le\exp\bigl(Ne^{-l/\xi}\bigr)-1
        \label{eq:hastings_percolation}
    \end{align}
    \cite[eq.~(perc2), lines 1227--1254]{Hastings2006Solving},
    \cite{gap:hastings05_percolation_cluster_size}.
\end{lemma}
\begin{proof}
    \leanok
    \uses{thm:hastings_percolation_exp, thm:hastings_percolation_cluster_sum}
    Combine (\ref{eq:hastings_percolation_exp}) and
    (\ref{eq:hastings_percolation_cluster_sum}) with
    $(xy)^{l+1}/(1-xy)\le(xy)^l=e^{-l/\xi}$ for $0<xy\le1/2$.
\end{proof}

\begin{lemma}[The percolation bound for connected clusters]
    \label{thm:hastings_percolation_connected}
    \lean{ClusterExpansion.sum_nonempty_powerset_prod_le_percolation_connected}
    \leanok
    \uses{thm:hastings_percolation, thm:hastings_connected_set_count}
    Let $G$ be a simple graph on $N$ vertices of maximum degree $\Delta$,
    $x=\max(1,\Delta^2)$, and $T$ be a family of connected vertex sets $C$ with
    $|C|\ge l+1$ and weights $0\le a_C\le y^{|C|}$. If $0<xy\le1/2$, then
    (\ref{eq:hastings_percolation}) holds with $\xi=1/\ln(1/(xy))$.
\end{lemma}
\begin{proof}
    \leanok
    \uses{thm:hastings_percolation, thm:hastings_connected_set_count}
    Combine Lemma~\ref{thm:hastings_connected_set_count}
    ($|\{C\in T:|C|=k\}|\le Nx^k$) with Lemma~\ref{thm:hastings_percolation}.
\end{proof}

\begin{lemma}[A percolation temperature]
    \label{thm:hastings_percolation_temperature}
    \lean{ClusterExpansion.exists_percolation_temperature}
    \leanok
    For any $J,S,x\in\mathbb R$ there exists $\beta_0>0$ such that
    $x(e^{J\beta}-1)e^{SJ\beta}\le1/2$ for all $0\le\beta\le\beta_0$
    \cite[lines 1242--1245]{Hastings2006Solving}.
\end{lemma}
\begin{proof}
    \leanok
    The function $\beta\mapsto x(e^{J\beta}-1)e^{SJ\beta}$ is continuous and
    vanishes at $\beta=0$.
\end{proof}

\begin{definition}[Clusters of large diameter]
    \label{def:hastings_large_diameter}
    \lean{ClusterExpansion.IsLargeDiam}
    \leanok
    \uses{def:hastings_cluster_of_terms}
    For terms $H_k$ supported in a metric space of sites, a set $C$ of terms has
    \emph{large diameter} if the support of $H_C$ contains two sites at distance
    $>l_{\mathrm{proj}}$ \cite[lines 1104--1107]{Hastings2006Solving}.
\end{definition}

\begin{lemma}[Size and number of clusters of large diameter]
    \label{thm:hastings_large_diameter_count}
    \lean{ClusterExpansion.dist_le_of_walk,
        ClusterExpansion.dist_le_of_mem_biUnion_of_connected,
        ClusterExpansion.card_ge_of_isLargeDiam,
        ClusterExpansion.card_add_card_boundary_le}
    \leanok
    \uses{def:hastings_large_diameter}
    If every term support has diameter at most $2R$ ($R>0$) and overlaps at most
    $\Delta$ other terms, then any connected set of $n$ terms has support
    diameter at most $2Rn$, every large-diameter connected set $C$ satisfies
    $n_C\ge\lfloor l_{\mathrm{proj}}/(2R)\rfloor+1$, and
    $n_C+n_{\partial C}\le(\Delta+1)n_C$
    \cite[line 1235]{Hastings2006Solving},
    \cite{gap:hastings05_percolation_cluster_size}.
\end{lemma}
\begin{proof}
    \leanok
    Any two sites in the support of $C$ are connected by an overlap walk of at
    most $n$ supports of diameter $\le2R$, giving distance $\le2Rn$, and each
    $k\in\partial C$ overlaps some $c\in C$.
\end{proof}

\begin{lemma}[The percolation bound for clusters of large diameter]
    \label{thm:hastings_percolation_large_diameter}
    \lean{ClusterExpansion.norm_trace_mul_exp_sub_truncatedExp_le_isLargeDiam,
        ClusterExpansion.exp_neg_floor_div_le,
        ClusterExpansion.norm_trace_mul_exp_sub_truncatedExp_le_isLargeDiam_dist}
    \leanok
    \uses{thm:hastings_large_diameter_count, thm:hastings_truncation_trace_bound,
        thm:hastings_percolation_connected}
    In the setting of Lemma~\ref{thm:hastings_large_diameter_count}, let
    $H=\sum_kH_k$ have $N$ Hermitian terms with $\|H_k\|\le J$, $\beta\ge0$,
    $x=\max(1,\Delta^2)$, $y=(e^{J\beta}-1)e^{(\Delta+1)J\beta}$, and
    $0<xy\le1/2$. With $Z=\operatorname{tr}e^{-\beta H}$,
    $l=\lfloor l_{\mathrm{proj}}/(2R)\rfloor$, $\xi=1/\ln(1/(xy))$, and
    $\xi_{\mathrm{perc}}=2R/\ln(1/(xy))$, for every $O$,
    \begin{align}
        \bigl|\operatorname{tr}\bigl(O(e^{-\beta H}-\rho(\beta,l_{\mathrm{proj}}))\bigr)\bigr|
         & \le\|O\|\bigl(\exp(Ne^{-l/\xi})-1\bigr)Z
        \notag                                                                               \\
         & \le\|O\|\bigl(\exp(N(xy)^{-1}e^{-l_{\mathrm{proj}}/\xi_{\mathrm{perc}}})-1\bigr)Z
        \label{eq:hastings_percolation_large_diameter}
    \end{align}
    \cite[eq.~(perc2), lines 1227--1254]{Hastings2006Solving},
    \cite{gap:hastings05_percolation_cluster_size,gap:hastings05_cluster_boundary}.
\end{lemma}
\begin{proof}
    \leanok
    \uses{thm:hastings_large_diameter_count, thm:hastings_truncation_trace_bound,
        thm:hastings_percolation_connected}
    Combine Lemma~\ref{thm:hastings_truncation_trace_bound} ($S=\Delta+1$),
    Lemma~\ref{thm:hastings_large_diameter_count}, and
    Lemma~\ref{thm:hastings_percolation_connected}, using
    $l\ge l_{\mathrm{proj}}/(2R)-1$.
\end{proof}

\section{The trace-norm bound and every temperature in Hastings' thermal expansion}
\label{sec:hastings_thermal_trace_norm}

% Hastings, arXiv:cond-mat/0508554v3, proj3new.tex, lines 1281--1301 (the
% trace-norm form) and 1303--1320 (eq. mplt, every temperature).

Two steps of \cite[Section~IV]{Hastings2006Solving} follow the high-temperature
cluster expansion. The first turns its bound on expectation values into a bound
in trace norm. The second extends the expansion to every temperature by taking an
even power of the high-temperature operator, eq.~(mplt) of
\cite[lines 1303--1317]{Hastings2006Solving}. This section proves the first step
and the positivity of the even power for general operators. The truncated
density matrix $\rho(\beta,l_{\mathrm{proj}})$ of the cluster expansion enters
only through the hypotheses. Its error bound at every temperature
\cite[lines 1322--1349]{Hastings2006Solving} is proved in
Section~\ref{sec:hastings_temperature_error}. The trace norm of a square matrix $D$ is
$\|D\|_1=\operatorname{Tr}\sqrt{D^\dagger D}$, and $\|O\|$ is the operator norm.

\begin{lemma}[The trace norm from expectation values]
    \label{thm:hastings_thermal_trace_norm_duality}
    \lean{Matrix.re_trace_sqrt_conjTranspose_mul_self_le_of_forall_norm_trace_mul_le}
    \leanok
    Let $D$ be a square complex matrix and $\epsilon$ a real number. If
    $|\operatorname{Tr}(OD)|\le\epsilon$ for every matrix $O$ with
    $\|O\|\le1$, then $\|D\|_1\le\epsilon$.
\end{lemma}
\begin{proof}\leanok
    The polar decomposition writes $D=U\sqrt{D^\dagger D}$ with $U$ unitary.
    Then $O=U^\dagger$ has $\|O\|\le1$ and
    $\operatorname{Tr}(OD)=\operatorname{Tr}\sqrt{D^\dagger D}=\|D\|_1$. For a
    Hermitian $D$ this $O$ may be taken to be the operator of signs of the
    eigenvalues of $D$ that \cite[lines 1291--1294]{Hastings2006Solving} uses;
    the polar decomposition removes the need for $D$ to be Hermitian.
\end{proof}

\begin{theorem}[The trace-norm form of the cluster-expansion bound]
    \label{thm:hastings_thermal_trace_norm}
    \lean{ClusterExpansion.re_trace_sqrt_sub_le_of_forall_norm_trace_mul_sub_le}
    \leanok
    Let $\rho$ and $\sigma$ be square complex matrices with
    $\operatorname{Re}\operatorname{Tr}\rho\ge0$, and let $\epsilon\ge0$. If
    \begin{align}
        |\operatorname{Tr}(O\rho)-\operatorname{Tr}(O\sigma)|
        \le\epsilon\,\|O\|\operatorname{Re}\operatorname{Tr}\rho
    \end{align}
    for every matrix $O$, then
    $\|\rho-\sigma\|_1\le\epsilon\operatorname{Re}\operatorname{Tr}\rho$. For the
    thermal state $\rho=\rho(\beta)$, which is positive with
    $\operatorname{Tr}\rho(\beta)=Z=\|\rho(\beta)\|_1$, and
    $\sigma=\rho(\beta,l_{\mathrm{proj}})$, the hypothesis is the bound of
    Lemma~\ref{thm:hastings_percolation_large_diameter} with
    $\epsilon=\exp[N(xy)^{-1}\exp(-l_{\mathrm{proj}}/\xi_{\mathrm{perc}})]-1$.
    The factor $(xy)^{-1}$ is the size-threshold correction to eq.~(perc2)
    \cite{gap:hastings05_percolation_cluster_size}. The conclusion is the
    corresponding trace-norm bound of \cite[lines 1281--1299]{Hastings2006Solving}.
\end{theorem}
\begin{proof}\leanok
    \uses{thm:hastings_thermal_trace_norm_duality}
    For $\|O\|\le1$ the hypothesis gives
    $|\operatorname{Tr}(O(\rho-\sigma))|\le\epsilon\operatorname{Re}\operatorname{Tr}\rho$,
    since $\epsilon\ge0$. Apply
    Lemma~\ref{thm:hastings_thermal_trace_norm_duality} to $D=\rho-\sigma$.
\end{proof}

\begin{lemma}[Even powers and the number of factors]
    \label{thm:hastings_even_powers}
    \lean{ClusterExpansion.posSemidef_pow_two_mul,
        ClusterExpansion.posSemidef_pow_of_even,
        ClusterExpansion.exists_even_div_le}
    \leanok
    Let $\rho$ be a Hermitian matrix and let $m$ be even. Then $\rho^m$ is
    positive semidefinite. For $\beta\ge0$ and $\beta_0>0$ there is an even
    $m\ge2$ with $\beta/m\le\beta_0$. This is the positivity part of eq.~(mplt)
    of \cite[lines 1303--1317]{Hastings2006Solving}, with positive
    semidefiniteness in place of the source's positive definiteness, which would
    need $\rho(\beta/m,l_{\mathrm{proj}})$ to be
    invertible~\cite{gap:hastings05_mplt_positivity}.
\end{lemma}
\begin{proof}\leanok
    For $m=2k$, $\rho^m=(\rho^k)^\dagger\rho^k$ since $\rho^k$ is Hermitian, and
    this is positive semidefinite. For the number of factors take
    $m=2(\lceil\beta/\beta_0\rceil+1)$.
\end{proof}

\begin{lemma}[Positivity of the truncated operator at every temperature]
    \label{thm:hastings_truncated_positivity}
    \lean{ClusterExpansion.isHermitian_clusterPart,
        ClusterExpansion.isHermitian_truncatedExp,
        ClusterExpansion.posSemidef_truncatedExp_pow}
    \leanok
    \uses{def:hastings_truncated_expansion, def:hastings_cluster_part}
    Let $H=\sum_kH_k$ with Hermitian terms $H_k$. Then every cluster part
    $\rho(\beta,C)$ and the truncated operator $\rho(\beta,l)$ are Hermitian, and
    for every even $m$ the power $\rho(\beta,l)^m$ is positive semidefinite. With
    $\beta/m$ in place of $\beta$, this is the positivity of
    $\rho(\beta,l_{\mathrm{proj}})=\rho(\beta/m,l_{\mathrm{proj}})^m$ in
    eq.~(mplt) of \cite[lines 1303--1314]{Hastings2006Solving}, with positive
    semidefiniteness in place of positive
    definiteness~\cite{gap:hastings05_mplt_positivity}.
\end{lemma}
\begin{proof}\leanok
    \uses{thm:hastings_truncation_inclusion_exclusion,
        thm:hastings_even_powers}
    Each $\rho(\beta,C)$ is a sum of the exponentials $e^{-\beta(H-H_B-H_D)}$ of
    Hermitian operators with coefficients $\pm1$, hence Hermitian. By
    Lemma~\ref{thm:hastings_truncation_inclusion_exclusion}, $\rho(\beta,l)$ is a
    sum of the $\rho(\beta,U)$ with coefficients $\pm1$, hence Hermitian, and
    Lemma~\ref{thm:hastings_even_powers} gives the positivity of
    its even powers.
\end{proof}

\begin{lemma}[The least sufficient even number of factors]
    \label{thm:hastings_even_temperature_factors}
    \lean{ClusterExpansion.evenTemperatureFactors,
        ClusterExpansion.evenTemperatureFactors_spec,
        ClusterExpansion.evenTemperatureFactors_le,
        ClusterExpansion.exp_evenTemperatureFactors_mul_sub_one_le}
    \leanok
    For $\beta>0$ and $\beta_0>0$, put
    $n=2\lceil\beta/(2\beta_0)\rceil$. Then $n\ge2$ is the least
    positive even integer with $\beta/n\le\beta_0$, and
    \begin{align*}
        \beta/\beta_0\le n<\beta/\beta_0+2.
    \end{align*}
    Consequently, for every $A\ge0$,
    $e^{nA}-1\le e^{(\beta/\beta_0+2)A}-1$.
    This supplies the exact rounding bound for the choice in
    \cite[lines 1307--1312]{Hastings2006Solving} and its corrected
    error estimate\cite{gap:hastings05_temperature_factor_count}.
\end{lemma}
\begin{proof}\leanok
    The ceiling lies between $\beta/(2\beta_0)$ and
    $\beta/(2\beta_0)+1$. Multiplication by two gives the displayed
    interval and $\beta/n\le\beta_0$. Any sufficient positive even
    integer is $2k$ with $k\ge\beta/(2\beta_0)$, hence is at least $n$.
    The last inequality follows from monotonicity of the exponential.
\end{proof}

\section{The error of the truncated expansion at every temperature}
\label{sec:hastings_temperature_error}

At an arbitrary inverse temperature $\beta$, the truncated operator is used in
$n$ factors: Hastings replaces $e^{-\beta H}=(e^{-\beta H/n})^n$ by the product
$\rho(\beta/n,l)^n$ of $n$ truncated factors, each at the inverse temperature
$\beta/n$ \cite[lines 1303--1314]{Hastings2006Solving}. The error is the sum of
the operators $\rho_m$ of the terms with at least $m$ large clusters across the
$n$ factors \cite[lines 1334--1341]{Hastings2006Solving}, and it is estimated
``following the same arguments as before''
\cite[lines 1342--1347]{Hastings2006Solving}. This section carries out that
estimate. Throughout, $y=(e^{J\beta/n}-1)e^{SJ\beta/n}$ and
$Z=\operatorname{tr}e^{-\beta H}$.

\begin{lemma}[The error of the product of truncated factors]
    \label{thm:hastings_temperature_error}
    \lean{ClusterExpansion.sum_prod_le_exp_of_subset_powerset,
        ClusterExpansion.norm_trace_mul_exp_sub_truncatedExp_pow_le}
    \leanok
    \uses{thm:hastings_truncation_inclusion_exclusion,
        thm:hastings_cluster_part_product, thm:hastings_percolation_exp}
    Let $H=\sum_kH_k$ be Hermitian with Hermitian terms of norm at most $J$, let
    $\beta\ge0$ and $n\ge1$, and let every large connected set $C$ have
    $n_C+n_{\partial C}\le Sn_C$. Then for every operator $O$
    \begin{align}
        \bigl|\operatorname{tr}\bigl(O(e^{-\beta H}-\rho(\beta/n,l)^n)\bigr)\bigr|
        \le\|O\|\Bigl(\exp\Bigl(n\sum_Cy^{n_C}\Bigr)-1\Bigr)Z,
        \label{eq:hastings_temperature_error}
    \end{align}
    summed over the large connected sets $C$.
\end{lemma}

\begin{proof}
    \leanok
    \uses{thm:hastings_truncation_inclusion_exclusion,
        thm:hastings_words_inclusion_exclusion, thm:hastings_cluster_part_product,
        thm:hastings_percolation_exp, thm:hastings_clusters_meet}
    By Lemma~\ref{thm:hastings_truncation_inclusion_exclusion} at the inverse
    temperature $\beta/n$, each factor is $\sum_U(-1)^{|U|}\rho(\beta/n,U)$ over
    the admissible families $U$ of large connected sets, and the empty family
    gives $e^{-\beta H/n}$. Expanding the $n$-th power gives a sum over the
    sequences $(U_0,\dots,U_{n-1})$ of admissible families of the ordered
    products $\prod_j(-1)^{|U_j|}\rho(\beta/n,U_j)$, and the sequence of empty
    families gives $e^{-\beta H}$. The members of an admissible family are
    pairwise disjoint (Lemma~\ref{thm:hastings_clusters_meet}), so the union
    of $U_j$ has $\sum_{C\in U_j}n_C$ terms and satisfies the hypothesis of
    Lemma~\ref{thm:hastings_cluster_part_product}, which bounds the term of
    every other sequence by $\|O\|Z\prod_jf(U_j)$ with
    $f(U)=\prod_{C\in U}y^{n_C}$. Summing over the sequences other than the
    empty one gives $\|O\|Z\bigl((\sum_Uf(U))^n-1\bigr)$. The nonempty admissible
    families are among the nonempty sets of large connected sets, so
    Lemma~\ref{thm:hastings_percolation_exp} gives
    $\sum_Uf(U)\le\exp(\sum_Cy^{n_C})$, and the $n$-th power of this is the
    exponential in~(\ref{eq:hastings_temperature_error}).
\end{proof}

Grouping the sequences by their total number $m$ of large clusters recovers the
bound $(1/m!)(n\sum_Cy^{n_C})^m$ on the operator $\rho_m$
of~\cite[lines 1334--1341]{Hastings2006Solving}.

\begin{theorem}[The error at every temperature, percolation form]
    \label{thm:hastings_temperature_error_percolation}
    \lean{ClusterExpansion.norm_trace_mul_exp_sub_truncatedExp_pow_le_percolation}
    \leanok
    \uses{thm:hastings_temperature_error, thm:hastings_percolation_cluster_sum,
        thm:hastings_percolation}
    In the setting of Lemma~\ref{thm:hastings_temperature_error}, let every
    large connected set have at least $l+1$ terms, let at most $Nx^k$ large
    connected sets have $k$ terms, for every $k$, where $N,x\ge0$, and let
    $0<xy\le1/2$. Put $\xi=1/\ln(1/(xy))$. Then for every operator $O$
    \begin{align*}
        \bigl|\operatorname{tr}\bigl(O(e^{-\beta H}-\rho(\beta/n,l)^n)\bigr)\bigr|
        \le\|O\|\bigl(\exp(nNe^{-l/\xi})-1\bigr)Z.
    \end{align*}
    The source prints the bound
    $\{\exp[(\beta/\beta_0)Ne^{-l/\xi}]-1\}\|O\|$ on
    $Z^{-1}|\operatorname{tr}O\rho(\beta)-\operatorname{tr}O\rho(\beta,l)|$
    \cite[lines 1342--1347]{Hastings2006Solving}, with $\beta/\beta_0$ in place of
    the number $n$ of factors. Its $n$, the smallest even integer with
    $\beta/n\le\beta_0$, satisfies $\beta/\beta_0\le n<\beta/\beta_0+2$, and its bound
    on $\rho_m$ gives the exponent with $n$, which is the bound proved here
    \cite{gap:hastings05_temperature_factor_count}. The threshold $l$ counts the
    terms of a large connected set rather than its diameter, as in
    Lemma~\ref{thm:hastings_percolation}
    \cite{gap:hastings05_percolation_cluster_size}.
\end{theorem}

\begin{proof}
    \leanok
    \uses{thm:hastings_temperature_error, thm:hastings_percolation_cluster_sum,
        thm:hastings_percolation}
    Lemma~\ref{thm:hastings_percolation_cluster_sum} bounds
    $\sum_Cy^{n_C}$ by $N(xy)^{l+1}/(1-xy)$, which is at most
    $N(xy)^l=Ne^{-l/\xi}$ since $xy\le1/2$, as in
    Lemma~\ref{thm:hastings_percolation}. Insert this into
    Lemma~\ref{thm:hastings_temperature_error}.
\end{proof}

\begin{corollary}[The corrected temperature-ratio bound]
    \label{cor:hastings_temperature_error_add_two}
    \lean{ClusterExpansion.norm_trace_mul_exp_sub_truncatedExp_pow_le_percolation_add_two}
    \leanok
    \uses{thm:hastings_temperature_error_percolation,
        thm:hastings_even_temperature_factors}
    In Theorem~\ref{thm:hastings_temperature_error_percolation}, take
    $\beta>0$, $\beta_0>0$ and
    $n=2\lceil\beta/(2\beta_0)\rceil$. Then
    \begin{align*}
        \bigl|\operatorname{tr}\bigl(O(e^{-\beta H}-\rho(\beta/n,l)^n)\bigr)\bigr|
        \le\|O\|\bigl(\exp[(\beta/\beta_0+2)Ne^{-l/\xi}]-1\bigr)Z.
    \end{align*}
    This is the explicit correction to the printed temperature factor
    in \cite[lines 1342--1347]{Hastings2006Solving}
    recorded in \cite{gap:hastings05_temperature_factor_count}.
\end{corollary}
\begin{proof}\leanok
    Apply $n<\beta/\beta_0+2$ in the percolation bound; the other
    factors and $Ne^{-l/\xi}$ are nonnegative.
\end{proof}

\section{Cluster truncation on a chain of interaction labels}
\label{sec:interval_cluster_truncation}

\subsection{Required labels with an empty forbidden set}

For a finite set $U$ of required interaction labels, let
\begin{align}
    R_U^H(\beta)
     & =\sum_{D\subseteq U}(-1)^{|D|}
    e^{-\beta(H-H_D)},
     & H_D                            & =\sum_{i\in D}H_i.
    \label{eq:mandatory_cluster_part}
\end{align}
This is the cluster part~(\ref{eq:hastings_cluster_part}) with the forbidden
set equal to $\varnothing$. The normalization below is an auxiliary
estimate for the interaction-label expansion related to the argument
of~\cite[lines 1270--1291 and 1935--1968]{Arad2016RigorousRG}.
It is distinct from the printed spatial approximation estimate, whose
obstructions are recorded in~\cite{gap:alvv16_cluster_approximation}.

\begin{lemma}[Normalized norm of the empty-boundary cluster part]
    \label{thm:mandatory_cluster_part_norm}
    \lean{ClusterExpansion.norm_clusterPart_empty_boundary_le_of_lower_bound}
    \leanok
    Let $H$ and $H_i$, $i\in U$, be Hermitian complex matrices of the same
    size. Suppose $\|H_i\|\leq1$, $H\geq\lambda I$, and $\beta\geq0$.
    Then
    \begin{align}
        e^{\beta\lambda}\|R_U^H(\beta)\|
        \leq(e^\beta-1)^{|U|}.
        \label{eq:mandatory_cluster_part_norm}
    \end{align}
    The real number $\lambda$ is any supplied lower energy. The background
    $H$ need not be the sum of the selected terms.
\end{lemma}
\begin{proof}\leanok
    \uses{def:hastings_cluster_part}
    The defining finite sum~(\ref{eq:mandatory_cluster_part}) is the
    alternating sum to which the normalized mandatory-label estimate
    applies, with no forbidden labels.
\end{proof}

\begin{lemma}[The mandatory-label part as a sum over exact supports]
    \label{thm:mandatory_cluster_part_support_sum}
    \lean{ClusterExpansion.clusterPart_empty_boundary_eq_sum_clusterOp}
    \leanok
    Let $K$ be a finite label set, let $H_i$, $i\in K$, be complex square
    matrices, and put $H=\sum_{i\in K}H_i$. For $U\subseteq K$ and every
    real $\beta$, the operators $E_A$ of
    Definition~\ref{def:hastings_cluster_operator} satisfy
    \begin{align}
        R_U^H(\beta)=\sum_{\substack{A\subseteq K \\U\subseteq A}}E_A.
        \label{eq:mandatory_cluster_part_support_sum}
    \end{align}
\end{lemma}
\begin{proof}\leanok
    \uses{thm:hastings_cluster_part_words, def:hastings_cluster_operator}
    At each word length, partition the words containing $U$ according to
    their exact letter set $A$. The possible sets $A$ form a finite
    family. Summing their convergent word series gives
    (\ref{eq:mandatory_cluster_part_support_sum}). Word order and repeated
    occurrences of a label are retained in this partition.
\end{proof}


\subsection{A normalized tail for an interaction-label cutoff}

The cutoff in this subsection counts interaction labels. It differs from
both physical-site support and the unrestricted estimate discussed
in~\cite{gap:alvv16_cluster_approximation}. No positivity of the individual
interactions is imposed.

\begin{definition}[Intervals of interaction labels]
    \label{def:cluster_label_intervals}
    \lean{ClusterExpansion.pathWindows,
        ClusterExpansion.mem_pathWindows,
        ClusterExpansion.pathIntervals,
        ClusterExpansion.mem_pathIntervals}
    \leanok
    Give the labels $K=\{0,\ldots,N-1\}$ their path adjacency. For an integer
    $r\geq1$, let $\mathcal W_r$ be the family of intervals of exactly $r$
    labels, and let $\mathcal I_r$ be the family of nonempty intervals of
    at least $r$ labels.
\end{definition}

\begin{lemma}[Windows and connected components on a path]
    \label{thm:cluster_label_window_geometry}
    \lean{ClusterExpansion.ordConnected_of_pathGraph_connected,
        ClusterExpansion.pathGraph_connected_of_ordConnected,
        ClusterExpansion.eq_Icc_of_ordConnected,
        ClusterExpansion.exists_interval_subset_card,
        ClusterExpansion.noLargeCluster_pathGraph_iff,
        ClusterExpansion.pathWindows_exists_marker,
        ClusterExpansion.card_pathIntervals_le,
        ClusterExpansion.clusters_biUnion_pathWindows_subset}
    \leanok
    \uses{def:cluster_label_intervals, def:hastings_cluster_of_terms}
    A nonempty support is connected in the label path if and only if it is
    an interval. Every interval with at least $r$ labels contains a member
    of $\mathcal W_r$. Hence a support has no component of size at least
    $r$ if and only if it contains no member of $\mathcal W_r$.
    The left endpoints of the members of $\mathcal W_r$ are distinct and
    belong to their respective windows. Moreover, $|\mathcal I_r|\leq N^2$.
    Every component of a union of members of $\mathcal W_r$ belongs to
    $\mathcal I_r$.
\end{lemma}
\begin{proof}\leanok
    \uses{thm:hastings_clusters_meet}
    A walk cannot cross a missing label: adjacent labels differ by one.
    Conversely, the induced path on an interval is the Hasse graph of its
    inherited order and is connected. A nonempty interval is determined
    by its minimum and maximum. Its size is the difference of these
    endpoints plus one, which gives a length-$r$ subinterval and the
    injectivity of the left endpoint for fixed length. The pair of
    endpoints gives the bound $N^2$.

    If a window meets a component of a support containing it, connectedness
    and avoidance of the component boundary force the window to lie in
    that component. This proves both the cutoff characterization and
    the assertion about unions of windows.
\end{proof}

\begin{lemma}[Inclusion and exclusion over mandatory windows]
    \label{thm:cluster_window_inclusion_exclusion}
    \lean{ClusterExpansion.sum_clusterOp_bad_windows_eq,
        ClusterExpansion.exp_sub_truncatedExp_eq_sum_clusterPart_windows,
        Finset.norm_sum_powerset_biUnion_sign_le_two_pow}
    \leanok
    \uses{thm:mandatory_cluster_part_support_sum}
    Let $K$ be finite, let $H_i$, $i\in K$, be complex square matrices,
    and put $H=\sum_iH_i$. Let $\mathcal W$ be a finite family of subsets
    of $K$. For every real $\beta$,
    \begin{align}
        \sum_{\substack{A\subseteq K                   \\
                \exists I\in\mathcal W:\ I\subseteq A}}E_A
         & =\sum_{\varnothing\ne F\subseteq\mathcal W}
        (-1)^{|F|+1}R_{\bigcup F}^H(\beta).
        \label{eq:cluster_window_inclusion_exclusion}
    \end{align}
    If the retained supports of a truncated expansion are exactly those
    avoiding every window, its error $e^{-\beta H}-M$ equals
    (\ref{eq:cluster_window_inclusion_exclusion}).
    If each member $I$ of $\mathcal W$ has a distinct selected label
    belonging to $I$, then the grouped coefficient
    \begin{align}
        c_U=\sum_{\substack{\varnothing\ne F\subseteq\mathcal W \\
                \bigcup F=U}}(-1)^{|F|+1}
        ,\qquad |c_U|\leq2^{|U|}.
        \label{eq:cluster_window_coefficient}
    \end{align}
\end{lemma}
\begin{proof}\leanok
    \uses{thm:mandatory_cluster_part_support_sum}
    Apply finite inclusion and exclusion to the events $I\subseteq A$,
    then sum over $A$. For each $F$, the intersection of these events is
    $\bigcup F\subseteq A$, whose exact-support sum is
    $R_{\bigcup F}^H(\beta)$.
    A contributing family with union $U$ consists of windows contained in
    $U$. Their selected labels are distinct elements of $U$, so there
    are at most $|U|$ such windows and at most $2^{|U|}$ contributing
    families. Each sign has absolute value one.
\end{proof}

\begin{lemma}[Weighted unions of windows]
    \label{thm:cluster_label_window_weights}
    \lean{ClusterExpansion.sum_pow_card_unions_pathWindows_le}
    \leanok
    \uses{thm:cluster_label_window_geometry}
    For $0\leq b\leq1$, put
    $\mathcal U_r=\{\bigcup F:\varnothing\ne F\subseteq\mathcal W_r\}$.
    Then
    \begin{align}
        \sum_{U\in\mathcal U_r}b^{|U|}
        \leq\exp(N^2b^r)-1.
        \label{eq:cluster_label_window_weights}
    \end{align}
\end{lemma}
\begin{proof}\leanok
    \uses{thm:cluster_label_window_geometry,
        thm:hastings_truncation_cluster_products}
    The component family determines $U$, its members are disjoint,
    and each belongs to $\mathcal I_r$. Thus
    \begin{align}
        \sum_{U\in\mathcal U_r}b^{|U|}
         & \leq\sum_{\varnothing\ne F\subseteq\mathcal I_r}
        \prod_{C\in F}b^{|C|}
        =\prod_{C\in\mathcal I_r}(1+b^{|C|})-1
        \leq\exp\Bigl(\sum_{C\in\mathcal I_r}b^{|C|}\Bigr)-1
        \leq\exp(N^2b^r)-1.
        \label{eq:cluster_label_weight_product}
    \end{align}
    The last inequality uses $|C|\geq r$ and
    $|\mathcal I_r|\leq N^2$.
\end{proof}

\begin{theorem}[Normalized interaction-label tail on a path]
    \label{thm:normalized_cluster_label_tail}
    \lean{ClusterExpansion.norm_exp_sub_truncatedExp_pathGraph_le}
    \leanok
    \uses{def:hastings_truncated_expansion, thm:mandatory_cluster_part_norm,
        thm:cluster_label_window_geometry, thm:cluster_window_inclusion_exclusion,
        thm:cluster_label_window_weights}
    Let $H=\sum_{i=0}^{N-1}H_i$, where the $H_i$ are Hermitian matrices
    with $\|H_i\|\leq1$. Suppose $H\geq\lambda I$, $r\geq1$,
    $\beta\geq0$, and $b=2(e^\beta-1)\leq1$.
    Let $M_r$ retain exactly the words whose label supports have no
    path component of size at least $r$. Then
    \begin{align}
        e^{\beta\lambda}\|e^{-\beta H}-M_r\|
        \leq\exp(N^2b^r)-1.
        \label{eq:normalized_cluster_label_tail}
    \end{align}
    The number $\lambda$ is any lower energy. No commutation of the terms
    is required for this estimate.
\end{theorem}
\begin{proof}\leanok
    \uses{thm:mandatory_cluster_part_norm, thm:cluster_label_window_geometry,
        thm:cluster_window_inclusion_exclusion, thm:cluster_label_window_weights}
    The cutoff characterization and inclusion and exclusion group the
    discarded operator by required label union.
    The left endpoints provide the distinct selected labels for the
    coefficient bound. Hence
    \begin{align}
        e^{\beta\lambda}\|e^{-\beta H}-M_r\|
         & \leq\sum_{U\in\mathcal U_r}|c_U|(e^\beta-1)^{|U|}
        \leq\sum_{U\in\mathcal U_r}b^{|U|}
        \leq\exp(N^2b^r)-1.
        \label{eq:normalized_cluster_label_tail_proof}
    \end{align}
    This changes both the normalization and the cutoff of the printed
    estimate in~\cite[lines 1270--1291 and 1935--1968]{Arad2016RigorousRG}.
    It does not assert that estimate or a matrix product representation.
\end{proof}


\subsection{Physical factors at a cut of the label path}

Let $K$ be a finite interaction-label set with graph $G$, let $H_i$ be
complex square matrices, and let $E_A$ be their exact-label operators.
For a fixed criterion of large components, write
\begin{align}
    F(B)=\sum_{\substack{A\subseteq B \\
            A\text{ has no large component}}}E_A.
    \label{eq:retained_allowed_labels}
\end{align}
Every operator in (\ref{eq:retained_allowed_labels}) is a finite sum of
the existing exact-label coefficients.

\begin{lemma}[Clusters of separated supports]
    \label{thm:clusters_separated_supports}
    \lean{ClusterExpansion.isCluster_union_iff_of_no_adj,
        ClusterExpansion.noLargeCluster_union_iff_of_no_adj,
        ClusterExpansion.isCluster_self_of_connected,
        ClusterExpansion.noLargeCluster_iff_not_large_of_connected}
    \leanok
    \uses{def:hastings_cluster_of_terms}
    Suppose no vertex of $B$ is adjacent to a vertex of $D$.
    A set $C$ is a component of $B\cup D$ if and only if it is a
    component of $B$ or of $D$. Hence $B\cup D$ has no large component
    if and only if neither $B$ nor $D$ has a large component.
    A connected nonempty set $C$ is its own component and has no
    large component if and only if $C$ itself is not large.
\end{lemma}
\begin{proof}\leanok
    \uses{thm:hastings_clusters_meet}
    The two supports have no edges between them. A connected component
    meeting one side therefore stays on that side. Conversely, a
    component of one side has no adjacent vertex in the other side.
    The assertions about a connected set follow from uniqueness of
    a component meeting it.
\end{proof}

\begin{lemma}[Finite retained sums on separated label sets]
    \label{thm:retained_separated_sum_product}
    \lean{ClusterExpansion.sum_retained_clusterOp_union_eq_mul}
    \leanok
    \uses{thm:clusters_separated_supports, def:hastings_cluster_operator}
    Suppose nonadjacent labels correspond to commuting matrices.
    If $B,D\subseteq K$ are disjoint and no edge joins them, then
    $F(B\cup D)=F(B)F(D)$.
\end{lemma}
\begin{proof}\leanok
    \uses{thm:clusters_separated_supports,
        thm:hastings_truncation_cluster_products}
    The maps $A\mapsto(A\cap B,A\cap D)$ and
    $(A_1,A_2)\mapsto A_1\cup A_2$ are inverse on the retained
    supports. The selector separates over the two parts, and
    $E_{A_1\cup A_2}=E_{A_1}E_{A_2}$. Distributing the two finite
    sums gives the identity.
\end{proof}

\begin{lemma}[Grouping by the component containing a label]
    \label{thm:retained_crossing_component_sum}
    \lean{ClusterExpansion.sum_retained_clusterOp_containing_eq}
    \leanok
    \uses{thm:clusters_separated_supports,
        thm:retained_separated_sum_product}
    Suppose nonadjacent labels correspond to commuting matrices,
    and fix $p\in K$. Let $\partial C$ be the graph boundary of $C$.
    Then
    \begin{align}
        \sum_{\substack{A\subseteq K  \\A\text{ has no large component}\\p\in A}}E_A
        =\sum_{\substack{C\subseteq K \\p\in C,\ C\text{ connected}\\C\text{ not large}}}
        E_C F\bigl(K\setminus(C\cup\partial C)\bigr).
        \label{eq:retained_crossing_component_sum}
    \end{align}
\end{lemma}
\begin{proof}\leanok
    \uses{thm:hastings_clusters_meet, thm:clusters_separated_supports}
    Choose the unique component $C$ of $A$ containing $p$ and put
    $B=A\setminus C$. Then $B$ avoids both $C$ and $\partial C$.
    Conversely, any retained $B\subseteq K\setminus(C\cup\partial C)$
    gives a retained support $C\cup B$ whose component containing $p$
    is exactly $C$. These maps are inverse, and the exact-label
    operator factors as $E_{C\cup B}=E_CE_B$.
\end{proof}

\begin{lemma}[The two exterior sides of an interval]
    \label{thm:cluster_path_exterior_sides}
    \lean{ClusterExpansion.pathGraph_exterior_eq_sides,
        ClusterExpansion.pathGraph_sides_separated,
        ClusterExpansion.pathGraph_compl_singleton_eq_sides,
        ClusterExpansion.pathGraph_point_sides_separated}
    \leanok
    \uses{thm:cluster_label_window_geometry}
    Give $K=\{0,\ldots,N-1\}$ its path adjacency. For a nonempty
    connected $C$, put
    \begin{align}
        L(C) & =\{j\in K:j+1<c\text{ for every }c\in C\},
             & R(C)                                       & =\{j\in K:c+1<j\text{ for every }c\in C\}.
        \label{eq:cluster_path_exterior_sides}
    \end{align}
    The exterior $K\setminus(C\cup\partial C)$ equals $L(C)\cup R(C)$,
    and these sets are disjoint and have no edge between them.
    Similarly, deleting a single label $p$ leaves the separated sets
    $L_p=\{j:j<p\}$ and $R_p=\{j:p<j\}$.
\end{lemma}
\begin{proof}\leanok
    \uses{thm:cluster_label_window_geometry}
    The set $C$ is an interval. A vertex outside it lies entirely
    on one side; avoiding its boundary increases the separation by
    one label. The displayed inequalities also prove disjointness
    and the absence of an edge. The assertion about a single deleted
    label follows by splitting its complement according to order.
\end{proof}

\begin{lemma}[Counting crossing label intervals]
    \label{thm:physical_cluster_crossing_count}
    \lean{ClusterExpansion.card_crossing_intervals_le}
    \leanok
    \uses{thm:cluster_label_window_geometry}
    For a fixed label $p$, there are at most $(r-1)^2$ connected
    label sets $C$ with $p\in C$ and $|C|<r$.
\end{lemma}
\begin{proof}\leanok
    \uses{thm:cluster_label_window_geometry}
    For $r\leq1$ there are none. Otherwise send $C=[a,b]$ to
    $(p-a,b-p)$. Each coordinate belongs to $\{0,\ldots,r-2\}$,
    and the two coordinates recover the endpoints and hence $C$.
\end{proof}

